\documentclass[10pt,a4paper]{amsart}
\usepackage[margin=19mm]{geometry}
\usepackage{amsmath,amssymb,mathtools}
\usepackage[scr=rsfs,cal=euler]{mathalfa}
\usepackage{xfrac}
\usepackage[unicode,hidelinks,bookmarks=false]{hyperref}
\newcommand{\ie}{i.e.\ }
\newcommand{\cf}{cf.\ }
\newcommand{\resp}{respectively\ }
\newcommand{\Subsection}{\subsection}
\providecommand{\nobreakdash}{\nobreak\hbox{-}}
\setlength{\emergencystretch}{2em}
\multlinegap0pt
\usepackage{algorithm}\usepackage{algpseudocode}
\usepackage{mathtools}
\usepackage{tikz}
\usepackage{amssymb}
\usepackage{xcolor}
\newtheorem{lemma}{Lemma}
\newtheorem{theorem}{Theorem}
\newcommand{\A}{\mathcal{A}}
\renewcommand{\L}{\mathcal{L}}
\newcommand{\M}{\mathcal{M}}
\renewcommand{\P}{\mathcal{P}}
\renewcommand{\S}{\mathbf{S}}
\title{A $q$-Polymatroid Framework for Information Leakage in Secure Linear Network Coding}
\author{Eimear Byrne, Johan Vester Dinesen, Ragnar Freij-Hollanti, Camilla Hollanti}
\date{Source: arXiv:2601.07567v1 (CC BY 4.0)}
\begin{document}
\maketitle

\noindent\emph{Source and licence.} A $q$-Polymatroid Framework for Information Leakage in Secure Linear Network Coding. Version arXiv:2601.07567v1 (CC BY 4.0), distributed by arXiv under the Creative Commons Attribution 4.0 International licence, http://creativecommons.org/licenses/by/4.0/, as recorded in the arXiv licence metadata for this exact version and shown on the abstract page https://arxiv.org/abs/2601.07567v1. Full licence notice: \texttt{licenses/arxiv-2601.07567-CC-BY-4.0.txt}.\par\smallskip

\noindent\textit{Editorial scope.} Counted unit: the article's theorem thm:brickelldavenport (source0.tex lines 804--805). The article's complete original proof of it is delivered with it (source0.tex lines 806--845, one \texttt{proof} environment of 40 lines). No mathematics is written, repaired or re-derived: the statement and the proof are the article's own lines, in the article's own notation, and no retained line is substituted or re-flowed.  Retained uncounted as the source-local context the proof needs: lines 776--778; lines 783--788.  Nothing was omitted from any retained span, so the package carries no omission record.  No retained line contains a citation, so the wrapper carries no bibliography and no bibliography entry.  No retained line contains a figure environment, an image-inclusion command or drawing code, so no image is delivered and the package declares no asset dependency.  Editorial formatting only, in addition: an AMS \texttt{amsart} preamble that restates the article's own declarations and macros used by the retained lines, namely the environments \texttt{lemma}, \texttt{theorem} on separate counters, together with the standard packages those lines need.  The retained environments print in this document as Lemma 1, Theorem 1 in order of appearance, whereas the article prints the same items with the numbering of its own full document.  The complete native source of the article as distributed in the arXiv e-print for this exact version is bundled unchanged as \texttt{sources/192514/source0.tex}.
\begin{lemma} Let $\S_{P_0,P}(\mathcal{M})$ be a perfect and connected $q$-polymatroid port such that $\rho(P_0)= \dim P_0 = 1$. For any $V\in \A$ there exists $W\in \A$ such that $V\cap W=0$ and $V+W\in \Gamma$, with the additional minimality property that if $W'<W$ then $V+W'\in \A$.
\label{lem:brickell1}
\end{lemma}

\begin{lemma}
Let $\S_{P_0,P}(\mathcal{M})$ be a perfect $q$-polymatroid port such that $\rho(P_0) = \dim P_0 =1$. Suppose $W\leq P$ and $y\in \P(P)$ such that both $W\in \A$ and $y+W\in \Gamma$. Then \label{lem:brickell2}
$$
    \rho(y\,|\, W) = 1 \quad \text{and} \quad \rho(y\,|\, W+P_0)=0.
$$
\end{lemma}

\begin{theorem}[$q$-Brickell-Davenport Theorem] Let $\S_{P_0,P}$ be an ideal, perfect and connected $q$-polymatroid port with $\rho(P_0) = \dim P_0 = 1$. Then $\M|_P$ is a $q$-matroid. \label{thm:brickelldavenport}
\end{theorem}

\begin{proof}
Towards a contradiction, suppose that there exists a subspace $V\leq P$ such that $\rho(V)$ is not a non-negative integer. Assume that with respect to inclusion, $V$ is a minimal subspace of $P$ with this property. Write $n= \lfloor \rho(V) \rfloor$, so $n < \rho(V) <n+1.$ For any $U\leq P$ and any $u\in \P(U)$ let $U_u$ denote a fixed hyperplane of $U$ that is a complement of $u$ in $U$, i.e. $U = U_u \oplus u$. Let $y\in \P(V)$, so $\rho(V_y)\in \mathbb{N}$, and as $\rho(y) = 1$ it follows that
\begin{align}
    \rho(V_y) \leq n <\rho(V) = \rho(V_y + y) \leq \rho(V_y) +1.\label{eq:brickell1}
\end{align}
This gives us that $n-1 <\rho(V_y) \leq n$ and since $\rho(V_y)$ is an integer we obtain $\rho(V_y) = n$. In particular,
\begin{align}
    0<\rho(y\,|\, V_y) <1. \label{eq:brickell2}
\end{align}
%
Since $\S_{P_0,P}$ is perfect, we have that the union of the elements of $\Gamma$ and $\A$ form a partition of $\L(P)$, and as such we consider the case that $V \in \Gamma$ and the case $V \in \A$.

Suppose first $V\in \Gamma$. For any $y\in \P(V)$ then $V_y\in \Gamma$ by Lemma \ref{lem:brickell2}, since if not, it would contradict $\eqref{eq:brickell2}$. Suppose now $V'\in \Gamma_{\min}$ such that $V'\leq V$ and let $y'\in \P(V')$. Then $V'_{y'}\in \A$ and Lemma \ref{lem:brickell2} yields $\rho(y'\,|\, V'_{y'}+P_0) = 0$, so
\[
    \rho(y'\,|\, V_{y'}) = \rho(V) - \rho(V_{y'}) = \rho(V+P_0) - \rho(V_{y'}) = \rho(y'\,|\, V_{y'}+P_0)=0,
\]
which contradicts \eqref{eq:brickell2} so $\rho(V)\in \mathbb{N}$.

Suppose now $V\in \A$. By Lemma \ref{lem:brickell1} there exists $W\in \A$ such that $V+W \in \Gamma$, $V\cap W = 0$ and if $W'<W$ then $V +W'\in \A$. Let $\dim W =\ell$ and choose $x_1,\dots,x_\ell \in \P(W)$ such that
\[
    W = x_1 \oplus x_2 \oplus \cdots \oplus x_\ell.
\]
For each $j\in [\ell]$ let $W_j = \oplus^\ell_{i=j}x_i$.  By Lemma \ref{lem:brickell2}, for any $x\in \P(W)$ we have $\rho(x\,|\, V+W_x) = 1$, which yields $\rho(x_j \,|\, V+W_{j+1})=1$ for all $j=1,\ldots,\ell-1$. We now have that
\begin{align*}
    \rho(V+W)   &= \rho(V+W_1)+\rho(x_1 \,|\, V+W_{x_1}) \\
                &= \rho(V+W_2) + \rho(x_2 \,|\, V+W_{x_2}) + \rho(x_1 \,|\, V+W_{x_1}) \\
                &\: \,\,\vdots \\
                &= \rho(V) + \sum^{\ell}_{i=1} \rho(x_i \,|\, V+W_i) \\
                &= \rho(V) + \dim W.
\end{align*}

Now, suppose $x\in \P(V)$. By Lemma \ref{lem:brickell2} we have that $V_x + W\in \Gamma$ for all $x\in \P(V)$ as $\rho(x\,|\, V_x + W)  \leq \rho(x\,|\, V_x) <1$. Furthermore, we have $\rho(V_x+W) = \rho(V) + \dim W$, as $\rho(V_x+W)\leq \rho(V) + \dim W$ holds trivially by submodularity, and $\rho(x\,|\, V_x + W) \leq \rho(x\,|\, V_x)$ implies $ \rho(V_x) + \dim W \leq \rho(V_x+W)$.


Choose $Z\in \Gamma$ that is minimal among those that contain $W$ and are contained in $V+W$. Thus, $W<Z \leq V+W$ and $V\cap Z >0$. Let $z\in \P(V\cap Z)$, so by Lemma \ref{lem:brickell2}, we have that $\rho(z\,|\, Z_z +P_0) = 0$, from which we obtain $\rho(z\,|\, V_z + W) = \rho(z\,|\, V_z +W+P_0) = 0.$ Therefore,
\[
    \rho(V) + \dim W = \rho(V+W) = \rho(V_z+W) + \rho(z\,|\, V_z+W) = \rho(V_z+W) = \rho(V_z) + \dim W,
\]
but this is a contradiction to \eqref{eq:brickell1}, which concludes the proof.
\end{proof}

\end{document}
